\thispagestyle{empty}
\section*{Anotácia}

\begin{minipage}[t]{1\columnwidth}%
Slovenská technická univerzita v Bratislave

FAKULTA INFORMATIKY A INFORMAČNÝCH TECHNOLÓGIÍ
\vskip 5pt
\begin{tabular}{@{}l@{\hspace{10pt}}>{\hyphenrules{nohyphenation}}p{9cm}@{}}

Študijný program: & \myStudyProgram \\

Autor: & \myName \\

Bakalárska práca: & \myTitle \\

Vedúci bakalárskej práce: & \mySupervisor \\
\end{tabular}
\\

\myDate%
\end{minipage}

\bigskip{}
\bigskip{}

\paragraph{}
V tejto diplomovej práci som sa najprv zameral na celú problematiku prezentovania, kde som sa napokon bližšie sústredil na verejné prezentácie. Najprv som sa venoval významu prezentácie ako takej a jej rozdeleniu. Následne som definoval pojem prezentačných schopností a rozdelil komunikáciu na verbálnu a neverbálnu časť. Práca taktiež obsahuje analýzu neverbálnej komunikácie a kladie dôraz na najpodstatnejšie negatívne javy. V práci je zahrnutá aj analýza vlastností hlasu a verbálnej komunikácie. V rámci analýzy je preskúmaná aj problematika časových radov, s dôrazom na segmentáciu a detekciu trendu. Pri porovnaní existujúcich riešení som sa sústredil na tri existujúce riešenia a porovnal som ich. Po analýze som navrhol komplexné riešenie, ktoré využíva mikroslužbovú architektúru pre dosiahnutie flexibilnosti riešenia. Pre používateľov som taktiež navrhol mobilnú aplikáciu pre nahrávanie videí a zobrazenie vizualizácií. Navrhnuté riešenie som napokon úspešne implementoval a otestoval.

\newpage{}
\thispagestyle{empty}

\newpage
\newpage
\thispagestyle{empty}
\mbox{}
\newpage
\thispagestyle{empty}

\section*{Annotation}

\begin{minipage}[t]{1\columnwidth}%
Slovak University of Technology Bratislava 

FACULTY OF INFORMATICS AND INFORMATION TECHNOLOGIES
\vskip 5pt
\begin{tabular}{@{}l@{\hspace{10pt}}>{\hyphenrules{nohyphenation}}p{9cm}@{}}

Degree course: & \myStudyProgramEng \\

Author: & \myName \\

Bachelor's thesis: & \myTitleEN \\

Supervisor: & \mySupervisor \\
\end{tabular}
\\

\myDateEng%
\end{minipage}

\bigskip{}
\bigskip{}

\paragraph{}
In this diploma thesis, I first focused on the overall issue of presenting, eventually narrowing my attention to public presentations. I began by addressing the significance of presentations as such and their classification. I then defined the concept of presentation skills and divided communication into verbal and non-verbal components. The thesis also includes an analysis of non-verbal communication and emphasizes the most significant negative phenomena. It further contains an analysis of voice characteristics and verbal communication. As part of the analysis, the topic of time series is also examined, with a focus on segmentation and trend detection. When comparing existing solutions, I focused on three current approaches and compared them. After the analysis, I proposed a comprehensive solution that uses a microservices architecture to achieve flexibility. I also designed a mobile application for users to record videos and display visualizations. Finally, I successfully implemented and tested the proposed solution.

\newpage
\thispagestyle{empty}
\mbox{}
\newpage